(1) 角动量守恒给出 \(\dot\varphi=\ell/(mr^2)\)。把切向动能并入径向势，严格能量积分为 \(E=\tfrac12m\dot r^2+U_{\rm eff}(r)\)，其中 \(U_{\rm eff}=(\ell^2/(2m)+g)/r^2\)。它随 \(r\) 严格递减，且从 \(+\infty\) 降至零，所以对每个 \(E>0\) 恰有一个 \(\dot r=0\) 的转折点。代入 \(\ell=mv_\infty b\) 得 \(r_{\min}^2=(\ell^2/(2m)+g)/E=b^2+2g/(mv_\infty^2)\)，因此 \(r_{\min}>b\)，排斥力确实阻止粒子更靠近中心。

(2) 令 \(u=1/r\)，实际径向力 \(F_r=-dU/dr=2g/r^3=2gu^3>0\)。Binet 方程 \(u''+u=-mF_r/(\ell^2u^2)\) 化为 \(u''+\kappa^2u=0\)，其中 \(\kappa^2=1+2mg/\ell^2>1\)，撇号表示对 \(\varphi\) 求导。选择最近接点 \(\varphi=0\)，有 \(u'(0)=0\)、\(u(0)=1/r_{\min}\)，从而 \(r(\varphi)=r_{\min}/\cos(\kappa\varphi)\)。物理轨道限于 \(|\varphi|<\pi/(2\kappa)\)，两端 \(r\to\infty\) 的渐近方向分别是 \(\varphi_\pm=\pm\pi/(2\kappa)\)，不能把极角变化误认成速度散射角。

(3) 入射与出射的轨道极角总变化为 \(\Delta\varphi=\pi/\kappa\)；自由直线应有 \(\Delta\varphi=\pi\)，故偏折角 \(\Theta=\pi-\Delta\varphi=\pi(1-1/\kappa)\)。记 \(a^2=2g/(mv_\infty^2)\)，于是 \(\Theta(b)=\pi[1-b/\sqrt{b^2+a^2}]\)。设 \(y=1-\Theta/\pi\in(0,1)\)，反解 \(b=a y/\sqrt{1-y^2}\)。轴对称中心力散射的环形面积元满足 \(d\sigma=2\pi b\,|db|\)、\(d\Omega=2\pi\sin\Theta\,d\Theta\)，所以 \(\displaystyle\frac{d\sigma}{d\Omega}=\frac{b}{\sin\Theta}\left|\frac{db}{d\Theta}\right|=\frac{a^2y}{\pi\sin\Theta(1-y^2)^2}\)。该式对有限 \(b>0\) 的单值偏折成立，且没有人为引入量子散射振幅。

(4) 取 \(a^2/b^2\ll1\)，二项式展开 \(1/\sqrt{1+a^2/b^2}=1-a^2/(2b^2)+O(a^4/b^4)\)，得到 \(\Theta\simeq\pi g/(mv_\infty^2b^2)\)。当 \(\Theta\to0^+\)，有 \(1-y^2\simeq2\Theta/\pi\)、\(\sin\Theta\simeq\Theta\)，因而 \(\displaystyle d\sigma/d\Omega\simeq\pi g/(2mv_\infty^2\Theta^3)\)。由于 \(g/(mv_\infty^2)\) 的量纲是长度平方，截面量纲正确；前向发散反映长程尾部对无限大碰撞参数的贡献，不能误解为有限探测角的无限计数。\(g\to0\) 时所有有限 \(b\) 的偏折趋零，\(b\to0^+\) 时偏折趋 \(\pi\)，恰与纯排斥中心势相符。
